\SourceLeaf{345}

\Needspace{8\baselineskip}\begin{center}
[III.]

\textbf{GEOMETRICAL RELATIONS BY MEANS OF COMPLEX NUMBERS.}

I.
\end{center}

$M$, $M'$ two points with complex numbers $m$, $m'$.

$P$ a variable point for which the azimuths $PM'-PM$ give a prescribed
difference $= \varphi$.

The locus of $P$ is a circle whose

\[
\text{centre } \frac{m\,e^{i\varphi}-m'\,e^{-i\varphi}}{2i\sin\varphi},
\]

\[
\text{radius } = \text{modulus of } \frac{m'-m}{2\sin\varphi}.
\]

\Needspace{8\baselineskip}\begin{center}
II.
\end{center}

$M$, $M'$, $m$, $m'$ as before.

$P$ a variable point whose distances from $M$, $M'$ are in a prescribed ratio $k:k'$.

The locus of $P$ is a circle whose

\[
\text{centre } \frac{m.k'k'-m'.kk}{k'k'-kk} = c,
\]

\[
\text{radius } = \text{modulus of } \frac{kk'}{k'k'-kk}\cdot(m'-m).
\]



\Needspace{7\baselineskip}\begin{center}
\textbf{REMARKS.}
\end{center}

The three preceding notes probably date from around 1840;
[I] is on a separate leaf, [II] and [III] are entered in a handbook.
A “handwritten note” already printed in vol.\,IV.\,p.\,397 also concerns
the circle. \hfill \textsc{Stäckel}.
